%! AP Statistics — Unit 4, Lesson 4.2 — Group Activity
\documentclass[10pt]{article}
\usepackage{apstats-article}
\usepackage{apstats-boxes}

\begin{document}

\pageheader{Unit 4, Lesson 4.2}{Group Activity: Genetics Simulation Lab}
\namepartnerperiod

\vspace{0.1in}

\begin{scenariobox}[Context]{sky}
\small
Two parents each carry one copy of a recessive gene for a certain condition. Each parent
passes one allele (D = dominant, r = recessive) to each child. A child has the condition
only if they receive two recessive alleles (rr). Work through the tier your teacher assigns.
\end{scenariobox}

\vspace{0.1in}

% ── Tier R ────────────────────────────────────────────────────────────────────
\begin{tcolorbox}[colback=goldbg, colframe=goldacc, boxrule=1pt, arc=2mm,
    left=4mm, right=4mm, top=3mm, bottom=3mm,
    title={\bfseries Tier R \enspace Remediate --- Run a Pre-Designed Simulation}]
\small

\textbf{Set-up:} Use a coin to represent each parent's allele contribution.
\begin{itemize}
  \item Heads (H) = dominant allele (D)
  \item Tails (T) = recessive allele (r)
\end{itemize}
One trial = flip two coins (one per parent). Record the two alleles the child receives.
The \textbf{event of interest}: child gets TT (= rr, affected).

\vspace{8pt}
\textbf{Run 20 trials.} Record your results in the table:

\vspace{6pt}
\renewcommand{\arraystretch}{1.5}
\begin{tabularx}{\linewidth}{@{} c c c c | c c c c | c c c c | c c c c | c c c c @{}}
1 & 2 & 3 & 4 & 5 & 6 & 7 & 8 & 9 & 10 & 11 & 12 & 13 & 14 & 15 & 16 & 17 & 18 & 19 & 20 \\
\hline
\blank{0.7cm} & \blank{0.7cm} & \blank{0.7cm} & \blank{0.7cm} &
\blank{0.7cm} & \blank{0.7cm} & \blank{0.7cm} & \blank{0.7cm} &
\blank{0.7cm} & \blank{0.7cm} & \blank{0.7cm} & \blank{0.7cm} &
\blank{0.7cm} & \blank{0.7cm} & \blank{0.7cm} & \blank{0.7cm} &
\blank{0.7cm} & \blank{0.7cm} & \blank{0.7cm} & \blank{0.7cm} \\
\end{tabularx}

\vspace{6pt}
(Write HH, HT, TH, or TT for each trial)

\begin{enumerate}[leftmargin=1.5em, itemsep=8pt]
  \item How many trials resulted in TT? \blank{2cm}
  \item What is your empirical probability? $\hat{p} = \blank{3cm}$
  \item Combine results with the class. Class total: \blank{2cm} TT out of \blank{2cm} trials.
        Class $\hat{p} = \blank{3cm}$
  \item The true probability is $\frac{1}{4} = 0.25$. How does the class estimate compare?\\
        \writeline
\end{enumerate}

\end{tcolorbox}

\vspace{0.1in}

% ── Tier A ────────────────────────────────────────────────────────────────────
\begin{tcolorbox}[colback=greenbg, colframe=greenacc, boxrule=1pt, arc=2mm,
    left=4mm, right=4mm, top=3mm, bottom=3mm,
    title={\bfseries Tier A \enspace Approaching Proficiency --- Full Simulation Design}]
\small

\textbf{Scenario:} A hospital screens newborns for a genetic condition. Each baby
independently has a 15\% chance of carrying at least one risk allele (based on population
genetics data). Use a random number table or random integers (00--99) to simulate 30
newborn screenings. Let 00--14 = ``carrier'' and 15--99 = ``non-carrier.''

\begin{enumerate}[leftmargin=1.5em, itemsep=8pt]
  \item Describe in your own words why assigning digits 00--14 correctly models a 15\%
        probability.\\
        \writeline\writeline
  \item Obtain a row of random two-digit numbers from your teacher (or generate them).
        Record 30 outcomes below:

        \vspace{4pt}
        \begin{tabularx}{\linewidth}{@{} *{15}{>{\centering\arraybackslash}X} @{}}
        \blank{0.5cm} & \blank{0.5cm} & \blank{0.5cm} & \blank{0.5cm} & \blank{0.5cm} &
        \blank{0.5cm} & \blank{0.5cm} & \blank{0.5cm} & \blank{0.5cm} & \blank{0.5cm} &
        \blank{0.5cm} & \blank{0.5cm} & \blank{0.5cm} & \blank{0.5cm} & \blank{0.5cm} \\
        \blank{0.5cm} & \blank{0.5cm} & \blank{0.5cm} & \blank{0.5cm} & \blank{0.5cm} &
        \blank{0.5cm} & \blank{0.5cm} & \blank{0.5cm} & \blank{0.5cm} & \blank{0.5cm} &
        \blank{0.5cm} & \blank{0.5cm} & \blank{0.5cm} & \blank{0.5cm} & \blank{0.5cm} \\
        \end{tabularx}

  \item Number of carriers: \blank{2cm} \quad Empirical probability: $\hat{p} = \blank{3cm}$
  \item How close is your estimate to the true probability of 0.15? \writeline
  \item If you repeated this simulation with 300 trials instead of 30, would you expect
        your estimate to be closer to 0.15 or farther? Explain using the Law of Large Numbers.\\
        \writeline\writeline
\end{enumerate}

\end{tcolorbox}

\vspace{0.1in}

% ── Tier E ────────────────────────────────────────────────────────────────────
\begin{tcolorbox}[colback=sky, colframe=navy, boxrule=1pt, arc=2mm,
    left=4mm, right=4mm, top=3mm, bottom=3mm,
    title={\bfseries Tier E \enspace Extension --- Design Your Own}]
\small

\textbf{Scenario:} A taste-test researcher wants to estimate the probability that a
randomly chosen person can correctly identify the brand of a cola drink in a blind
taste test. Without any real ability, each guess is equally likely to be correct or
incorrect.

\begin{enumerate}[leftmargin=1.5em, itemsep=8pt]
  \item Design a simulation to estimate the probability that a guesser with no real
        ability gets at least 7 of 10 guesses correct. Describe your simulation clearly:
        \begin{enumerate}[label=(\alph*), itemsep=4pt]
          \item Chance device: \writeline
          \item How values are assigned: \writeline
          \item What constitutes one trial: \writeline
          \item What the event of interest is: \writeline
        \end{enumerate}
  \item Run 25 trials and record outcomes:
        \vspace{20pt}
  \item Number of trials in which at least 7 of 10 were correct: \blank{2cm}\\
        Empirical probability: $\hat{p} = \blank{3cm}$
  \item The true probability of getting 7 or more correct by chance is approximately
        0.172. How does your estimate compare? What would happen with more trials?\\
        \writeline\writeline
\end{enumerate}

\end{tcolorbox}

\end{document}
